% ----------------------------------------------------------
% Seção 4 — Referências (NBR 6023)
% ----------------------------------------------------------
\chapter{Referências}\label{cap:referencias}

Referência é o conjunto padronizado de elementos que permite identificar um documento \cite{abnt6023}. No \LaTeX, as referências não são digitadas na lista final: cada documento é cadastrado uma vez no arquivo \texttt{referencias.bib}, e a lista é montada com as obras efetivamente citadas, em ordem alfabética e no formato da versão da ABNT escolhida.

\section{O arquivo .bib}\label{sec:arquivo-bib}

Cada entrada do arquivo \texttt{.bib} tem um tipo, uma chave e uma lista de campos:

\begin{minted}{bibtex}
@book{gil2017,
  author    = {Antonio Carlos Gil},
  title     = {Como elaborar projetos de pesquisa},
  edition   = {6},
  address   = {São Paulo},
  publisher = {Atlas},
  year      = {2017}
}
\end{minted}

A chave (\texttt{gil2017}) é o nome usado nos comandos de citação. Alguns cuidados:

\begin{alineas}
	\item vários autores são separados por \texttt{and}, nunca por vírgula ou ponto e vírgula;
	\item entidades como autor ficam entre chaves duplas, para não serem lidas como nome e sobrenome: \verb|{{ASSOCIAÇÃO BRASILEIRA DE NORMAS TÉCNICAS}}|, em maiúsculas para que a versão 2002 as mantenha na citação;
	\item subtítulos vão no campo \texttt{subtitle}, para que só o título receba destaque;
	\item siglas e nomes próprios em títulos em inglês ficam entre chaves, para preservar as maiúsculas: \verb|{MapReduce}|;
	\item documentos on-line levam a data de acesso em dois campos: \texttt{urlaccessdate} (versão 2002) e \texttt{urldate}, no formato AAAA-MM-DD (versões 2018 e 2023).
\end{alineas}

A maioria das bases de periódicos, catálogos de bibliotecas e o Google Acadêmico exportam referências em formato Bib\TeX. Revise sempre a entrada exportada: campos faltantes ou com maiúsculas indevidas são comuns.

\section{Tipos de documento}\label{sec:tipos-documento}

O \autoref{qua:tipos-bib} relaciona os tipos de documento mais comuns em trabalhos acadêmicos ao tipo de entrada no \texttt{.bib}, com um exemplo citado neste guia. As referências completas estão na lista de Referências ao final do documento.

\begin{quadro}[htbp]
	\caption{Tipos de documento e entradas no arquivo .bib}\label{qua:tipos-bib}
	\centering
	\small
	\begin{tabular}{p{0.27\textwidth}p{0.22\textwidth}p{0.37\textwidth}}
		\toprule
		\textbf{Documento} & \textbf{Entrada} & \textbf{Exemplo} \\
		\midrule
		Livro & \texttt{@book} & \cite{lamport1994} \\
		Capítulo de livro & \texttt{@incollection} & \cite{brooks1995} \\
		Artigo de periódico & \texttt{@article} & \cite{wilson2017} \\
		Trabalho em evento & \texttt{@inproceedings} & \cite{dean2004} \\
		Tese & \texttt{@phdthesis} & \cite{fielding2000} \\
		Dissertação & \texttt{@mastersthesis} & \cite{shannon1940} \\
		Trabalho de conclusão de curso & \texttt{@mastersthesis} com \texttt{type} & ver \autoref{sec:tcc} \\
		Norma técnica & \texttt{@book} & \cite{abnt14724} \\
		Legislação & \texttt{@misc} & \cite{brasil1998} \\
		Página da internet & \texttt{@misc} com \texttt{url} & \cite{semver} \\
		Programa de computador & \texttt{@manual} & \cite{biblatexabnt} \\
		Conjunto de dados & \texttt{@misc} com \texttt{doi} & \cite{fisher1936} \\
		\bottomrule
	\end{tabular}
	\fonte{Elaborado pelo próprio autor.}
\end{quadro}

\subsection{Livro e capítulo de livro}

O livro tem como elementos essenciais autor, título, edição (a partir da segunda), local, editora e ano, como em \citeonline{chacon2014}. No capítulo, o autor e o título do capítulo vêm antes da expressão ``In:'' e dos dados do livro. O capítulo de \citeonline{brooks1995} está na edição comemorativa de um livro do próprio autor; nesse caso, o campo \texttt{editor} repete o nome do autor.

\subsection{Artigo de periódico}

O artigo tem como elementos o título do periódico, volume, número, páginas e ano. Quando houver, o DOI é informado no campo \texttt{doi}, como em \citeonline{knuth1984literate}.

\subsection{Trabalho em evento}

O trabalho apresentado em evento é registrado com \texttt{@inproceedings}, com o título dos anais em \texttt{booktitle} e o nome do evento em \texttt{eventtitle}, como em \citeonline{dean2004}.

\subsection{Tese, dissertação e trabalho de conclusão de curso}\label{sec:tcc}

Teses usam \texttt{@phdthesis} e dissertações, \texttt{@mastersthesis}, com a instituição no campo \texttt{school}: \citeonline{fielding2000} e \citeonline{shannon1940}. Para um trabalho de conclusão de curso, use \texttt{@mastersthesis} e informe o tipo de trabalho no campo \texttt{type}. O modelo de entrada a seguir mostra os campos; substitua cada valor pelos dados do trabalho consultado:

\begin{minted}{bibtex}
@mastersthesis{sobrenome2025,
  author  = {Nome Sobrenome},
  title   = {Título do trabalho},
  type    = {Trabalho de Conclusão de Curso (Bacharelado em
             Sistemas de Informação)},
  school  = {Universidade Federal de Viçosa},
  address = {Rio Paranaíba},
  year    = {2025}
}
\end{minted}

\subsection{Norma técnica e legislação}

Normas técnicas têm a entidade como autor e o código da norma como título, com o nome da norma em \texttt{subtitle}, como em \citeonline{abnt6023}. Leis e outros atos normativos têm a jurisdição como autor, a epígrafe (tipo, número e data) como título e a ementa e os dados de publicação em \texttt{note}, como em \citeonline{brasil1998}.

\subsection{Documentos on-line, programas e conjuntos de dados}

Documentos consultados on-line incluem o endereço eletrônico e a data de acesso, como em \citeonline{semver}. Programas de computador e pacotes \LaTeX\ registram a versão utilizada, como em \citeonline{biblatex}. Conjuntos de dados indicam o repositório e o DOI, como em \citeonline{fisher1936}.

\section{O que muda entre as versões}\label{sec:mudancas-referencias}

Entre a NBR 6023 de 2002 \cite{abnt6023-2002} e a de 2018 \cite{abnt6023}, as mudanças mais visíveis na lista de referências são:

\begin{alineas}
	\item as expressões para local e editora não identificados passam de [S.l.] e [s.n.] para \textit{[S. l.]} e \textit{[s. n.]};
	\item o endereço eletrônico deixa de ser escrito entre os sinais \texttt{< >};
	\item o DOI passa a ser um elemento complementar previsto pela norma;
	\item as referências passam a ser separadas entre si por uma linha em branco de espaço simples.
\end{alineas}

Compile o guia nas três versões e compare as listas de referências para ver essas diferenças.

A edição de 2025 da NBR 6023 \cite{abnt6023-2025} alterou os elementos das referências de coleções e fascículos de periódicos, de artigos eletrônicos (com a inclusão da e-localização), de eventos no todo e de jurisprudência. O estilo \texttt{biblatex-abnt} segue a edição de 2018; ao referenciar esses tipos de documento, confira o resultado com a edição de 2025.
